\subsection{Описание метода}

\subsubsection{Начальные данные}

Рассмотрим систему нелинейных уравнений размерности $n$ в векторном виде:
\[
\vec{F}(\vec{x}) = \vec{0}, \quad
\vec{F}(\vec{x}) = \bigl(f_1(x_1,\dots,x_n),\dots,f_n(x_1,\dots,x_n)\bigr)^T.
\]

Для применения метода простых итераций исходную систему необходимо преобразовать к эквивалентному виду
\[
\vec{x} = \Phi(\vec{x}),
\]
где $\Phi: \mathbb{R}^n \to \mathbb{R}^n$ --- некоторая вектор-функция.  
Такое преобразование не~является единственным. Сходимость метода напрямую зависит от выбора $\Phi$.

\subsubsection{Условия сходимости}

Для сходимости итерационного процесса $\vec{x}^{(k+1)} = \Phi(\vec{x}^{(k)})$ исследуем невязку:
\[
  \vec{e}^{(k)} = \vec{x}^{(k)} - \vec{x}^*
\]

Рассмотрим класс функций, дифференцируемых в точке \( \vec{x}^* \). Разложим $\Phi(\vec{x})$ в ряд Тейлора в окрестности $\vec{x}^*$:
\[
\Phi(\vec{x}^{(k)}) = \Phi(\vec{x}^*) + J_\Phi(\vec{x}^*)(\vec{x}^{(k)} - \vec{x}^*) + O(\|\vec{x}^{(k)} - \vec{x}^*\|^2)
\]
где $J_\Phi(\vec{x}^*)$ — матрица Якоби. Учтём, что $\vec{x}^{(k+1)} = \Phi(\vec{x}^{(k)})$ и $\vec{x}^* = \Phi(\vec{x}^*)$:
\[
\begin{gathered}
  \vec{x}^{(k+1)}\approx\vec{x}^*+J_\Phi(\vec{x}^*)(\vec{x}^{(k)}-\vec{x}^*)\\
  \vec{e}^{(k+1)} \approx J_\Phi(\vec{x}^*) \vec{e}^{(k)}
\end{gathered}
\]

То есть условие сходимости сводится к требованию, чтобы спектральный радиус матрицы Якоби в точке решения был меньше единицы:
\[
q\equiv\rho(J_\Phi(\vec{x}^*)) < 1
\]

Из этого можем вывести следствие: если в некоторой окрестности решения существует такая согласованная матричная норма, что:
\[
\|J_\Phi(\vec{x})\| \le q < 1,
\]
то итерационная последовательность сходится к решению $\vec{x}^*$ \textit{линейно}:
\[
\|\vec{e}^{(k+1)}\| \le q \|\vec{e}^{(k)}\|
\]

\subsubsection{Критерий останова и его обоснование}

Для контроля точности итерационного процесса необходимо установить связь между расстоянием до точного решения $\|\vec{x}^{(k)} - \vec{x}^*\|$ и разностью двух последовательных приближений $\|\vec{x}^{(k)} - \vec{x}^{(k-1)}\|$.

Оценим расстояние от текущей итерации до корня как сумму всех последующих шагов через неравенство треугольника:
\[
\|\vec{x}^{(k)} - \vec{x}^*\| = \left\| \sum_{i=k}^{\infty} (\vec{x}^{(i)} - \vec{x}^{(i+1)}) \right\| \le \sum_{i=k}^{\infty} \|\vec{x}^{(i)} - \vec{x}^{(i+1)}\|
\]

Учтём, что каждый шаг уменьшается пропорционально $q$:
\[
\|\vec{x}^{(i)} - \vec{x}^{(i+1)}\| \le q^{i-k} \|\vec{x}^{(k)} - \vec{x}^{(k+1)}\|
\]
Подставляя это в сумму, получим убывающую геометрическую прогрессию:
\[
\|\vec{x}^{(k)} - \vec{x}^*\| \le \|\vec{x}^{(k)} - \vec{x}^{(k+1)}\| \cdot (1 + q + q^2 + \dots) = \frac{1}{1-q} \|\vec{x}^{(k)} - \vec{x}^{(k+1)}\|
\]

Воспользуемся поправкой на один шаг и получим финальную оценку:
\[
\|\vec{x}^{(k)} - \vec{x}^*\| \le \frac{q}{1-q} \|\vec{x}^{(k)} - \vec{x}^{(k-1)}\|
\]

Таким образом, для достижения заданной точности $\varepsilon$ расчёт необходимо вести до выполнения условия:
\[
\|\vec{x}^{(k)} - \vec{x}^*\| \le \varepsilon\implies\|\vec{x}^{(k)} - \vec{x}^{(k-1)}\| \le \frac{1-q}{q} \varepsilon
\]

Если метод сходится быстро, то коэффициент перед погрешностью в финальной оценке становится меньше или равен единице:
\[
  q \le 0.5 \implies \frac{q}{1-q} \le 1.
\]
В этом случае условие $\|\vec{x}^{(k)} - \vec{x}^{(k-1)}\| \le \varepsilon$ гарантирует, что истинная погрешность $\|\vec{x}^{(k)} - \vec{x}^*\|$ также не превышает $\varepsilon$.